%قبل از شروع مسئله ترجیحاً چیزی ننویسید. اما اگر خواستید بنویسید، در قالب لیستی (فقرات و شمارش) بنویسید.
%برای اضافه کردن نمره، بعدی از شروع‌مسئله، [x] را اضافه کنید و به جای x  نمره مورد نظر را قرار دهید. برای زیرمسئله‌ها هم به همین صورت است.
%\شروع‌مسئله[10]
\شروع‌مسئله
% صورت اصلی مسئله را در این‌جا بنویسید.
در این مسئله می‌خواهیم شیوه‌ای در مورد پیاده‌سازی پیمایش میان‌ترتیب درخت دودودیی جست‌وجو
\LTRfootnote{inorder}
را بررسی و تحلیل کنیم.


یکی از روش‌های پیاده‌سازی میان‌ترتیب، استفاده از دو تابع به نامی های $tree\_min$ و $inorder\_successor$ است. تابع اول یک ریشه ورودی گرفته و عنصر کمینه زیردرخت نظیر آن (چپ‌ترین گره) را بر می‌گرداند. تابع دوم یک گره را ورودی گرفته و گره بعدی در پیمایش میان‌ترتیب را پیدا می‌کند. پیاده‌سازی آن به شکل زیر است:

\begin{latin}
\begin{algorithm}[H]
	\begin{algorithmic}
	
		\Procedure{$inorder\_successor(x)$}{}
		\If{$x.right \neq \text{null}$}
		\Return  $tree\_min(x.right)$
		\EndIf
		\State $y = x.parent$
		\While {$y \neq \text{null} ~~ \text{\textbf{and}} ~~ x == y.right$}
		\State $x = y$
		\State $y = y.p$
		\EndWhile
		\Return $y$
		\EndProcedure
	\end{algorithmic}
\end{algorithm}	
	
\end{latin}

توجه الگوریتم بالا معادل الگوریتم صفحه $83$ اسلایدهاست.
%اگر مسئله چند قسمت داشت به این صورت عمل کنید:
\شروع{زیرمسائل}

\زیرمسئله
فرض کنید $T$ یک درخت دودویی جست‌وجو با عناصر متمایز باشد. نشان دهید که اگر $x$ زیردرخت راست نداشته باشد و آخرین عنصر در پیمایش میان‌ترتیب هم نباشد،‌ آن‌گاه عنصر بعدی در پیمایش میان‌ترتیب بعد از $x$، پایین‌ترین جدی از $x$ خواهد بود که فرزند سمت چپ آن هم جد $x$ باشد. به بیان دیگر، این سوال قصد اثبات درستی تابع $inorder\_successor$ را دارد. 
\زیرمسئله
یک راه برای چاپ پیمایش میان‌ترتیب درخت دودویی جست‌وجو، این است که ابتدا یک بار تابع $tree\_min$ را روی ریشه کل درخت صدا کرده و پس از آن، $n-1$ بار تابع $inorder\_successor$ را روی خروجی مرحله قبل صدا کنیم. این روش به وضوح درست است. مرتبه زمانی اجرای این الگوریتم را بدست آورید.
\پایان{زیرمسائل}
\solution{
	فرض کنید راسی که مسئله خواسته است را با $y$ نمایش دهیم. ابتدا نشان می‌دهیم که $y$ باید جد $x$ باشد. اگر این طور نباشد، فرض کنید $z$ اولین جد مشترک $x$ و $y$ باشد. با توجه به این که $x$ فرزند سمت راست هم ندارد طبق خاصیت درخت دودویی جست‌وجو باید داشته باشیم $x < z < y$. در این صورت $y$ نمی‌تواند عنصر  بلافاصله بعد از $x$ در پیمایش میان‌ترتیب باشد. در نتیجه فرض خلف باطل است و $y$ باید جد $x$ باشد.
	
	حال توجه کنید که $y.left$ باید جد $x$ باشد، زیرا اگر این طور نباشد، $y.right$ جد $x$ خواهد بود. اگر این طور باشد یعنی $x>y$ که باز هم نتیجه می‌دهد که $y$ نمی‌تواند عضو بعدی در پیمایش میان‌ترتیب باشد. در نهایت نشان می‌دهیم که $y$ باید پایین‌ترین جد $x$ باشد که فرزند چپ آن هم جد $x$ است. فرض کنید $y$ پایین ترین جد با این خاصیت نباشد. در آن صورت $z$ با این شرایط وجود دارد. این بدین معنی خواهد بود که $z$ خودش در زیردرخت سمت چپ $y$ است که یعنی $z<y$ که این هم با نحوه پیماش میان ترتیب و این که $y$ عضو بعدی در این پیمایش بوده است، تناقض خواهد داشت.
	
	برای تحلیل زمانی الگوریتم،‌ نشان‌می‌دهیم که این الگوریتم از
	$\Theta(n)$
	است. برای این کار نشان می‌دهیم که الگوریتم پیمایش میان‌ترتیب به شکل گفته شده، هر یال را دو بار پیمایش می‌کند و از آن جایی که تعداد یال ها یکی کمتر از راس‌هاست، الگوریتم در کل از $\Theta(n)$ خواهد بود.
	
	برای اثبات این موضوع یک راس را به نام $x$ در درخت دودویی جست‌وجو در نظر بگیرید. حال یال بین $x$ و $x.parent$ را در نظر بگیرید. این یالیک بار در هنگامی که $inorder\_successor$ روی $x.parent$ فراخوانی می‌شود پیموده شده و بار دیگر، زمانی که این تابع روی بزرگترین عنصر موجود در زیردرخت با ریشه $x$ صدا زده می‌شود، برای یافتن جد آن پیموده خواهد شد. از آن جایی که همه یال‌ها با این استدلال تنها دو بار پیموده می‌شوند و یک هزینه اولیه هم برای $tree\_min$ وجود دارد که تاثیری در مرتبه زمانی نهایی الگوریتم نمی‌گذارد. در نتیجه زمان اجرا از $O(n)$ خواهد بود و همچنین از آن‌جایی که اندازه خروجی $n$ است،  از $\Omega(n)$ هم هست. در نتیجه این الگوریتم از $\Theta(n)$ است.



}



%%%%%%%%%%%%%%%%%%%%%%%%%%%%%%%%%%%%%%%%%%%%%%%%%%%%%
% برخی دستورات مفید برای تمرین‌ها. %
%%%%%%%%%%%%%%%%%%%%%%%%%%%%%%%%%%%%%%%%%%%%%%%%%%%%%

% فرمول‌های نمونه:
%$f(x) = \Theta(x)$
%$T(x, y) \leq \log(x) + 2^y + \binom{2n}{n}$
% {\tt code\_function}


% لیست موارد بدون شماره‌گذاری:
%\شروع{فقرات}
%    \فقره مورد اول 
%        \شروع{فقرات}
%            \فقره مورد اول داخلی 
%                \شروع{فقرات}
%                    \فقره اولین مورد خیلی داخلی 
%                    \فقره دومین مورد خیلی داخلی 
%                \پایان{فقرات}
%            \فقره مورد دوم داخلی 
%        \پایان{فقرات}
%    \فقره مورد دوم 
%\پایان{فقرات}

% لیست موارد شماره‌گذاری شده:
%\شروع{شمارش}
%    \فقره یک مورد 
%    \فقره یک مورد دیگر 
%    \فقره باز هم یک مورد دیگر
%\پایان{شمارش}

% برای نوشتن یک فرمول چند خطی که یک شماره بخوره. برای حالت‌های دیگر به سایت زیر مراجعه کنید:
% https://www.overleaf.com/learn/latex/Aligning_equations_with_amsmath
%\begin{align} 
%    \begin{split}
%        (x+y)^3 &= (x+y)^2(x+y) \\
%                &= (x^2+2xy+y^2)(x+y) \\
%                &= (x^3+2x^2y+xy^2) + (x^2y+2xy^2+y^3) \\
%                &= x^3+3x^2y+3xy^2+y^3
%    \end{split}                                 
%\end{align}

% نوشتن ماتریس و آرایه دوبعدی
%\begin{align}
%    A = 
%    \begin{bmatrix}
%        A_{11} & A_{21} \\
%        A_{21} & A_{22}
%    \end{bmatrix}
%\end{align}

% اضافه کردن تصویر با شرح. برای اضافه کردن تصویر بدون شرح، فقط دستور درج‌تصویر را استفاده کنید.
%\شروع{شکل}[h!]
%\تنظیم‌ازوسط
%  \درج‌تصویر[پهنا=0.5\پهنای‌متن]{img.jpg}
%\شرح{شرح تصویر}
%\برچسب{تصویر:۱}
%\پایان{شکل}
%\پایان{زیرمسائل}

%کد زیر را ببینید.
%%\lstinputlisting{source_filename.py}
%\begin{latin}
%\begin{lstlisting}
%# Import the modules
%import sys
%import random
%
%ans = True
%
%while ans:
%question = raw_input("Ask the magic 8 ball a question: (press enter to quit) ")
%
%answers = random.randint(1,8)
%
%if question == "":
%sys.exit()
%
%elif answers == 1:
%print "It is certain"
%
%elif answers == 2:
%print "Outlook good"
%
%elif answers == 3:
%print "You may rely on it"
%
%elif answers == 4:
%print "Ask again later"
%
%elif answers == 5:
%print "Concentrate and ask again"
%
%elif answers == 6:
%print "Reply hazy, try again"
%
%elif answers == 7:
%print "My reply is no"
%
%elif answers == 8:
%print "My sources say no"
%\end{lstlisting}
%\end{latin}
